\documentclass[11pt]{article}
\usepackage[margin=1in]{geometry}
\usepackage{amsmath,amssymb,amsthm}
\usepackage{lmodern,microtype}
\usepackage[T1]{fontenc}
\usepackage{booktabs,array}
\usepackage[hidelinks]{hyperref}
\newtheorem{theorem}{Theorem}
\newtheorem{lemma}{Lemma}
\newcommand{\R}{\mathbb R}
\newcommand{\Z}{\mathbb Z}
\newcommand{\N}{\mathbb N}
\DeclareMathOperator{\tr}{tr}
\title{Disproof of Conjecture 00000000236\\[3pt]
\large A uniform gap for normalized determinants at odd orders}
\author{C0ldSmi1e}
\date{October 5, 2026}
\begin{document}
\maketitle

\begin{abstract}
Let $D(n)$ be the maximum of $|\det A|$ over all $n\times n$ matrices
with entries in $\{-1,1\}$. For every odd $n\geq 3$, we prove
$D(n)^2/n^n\leq \exp(-1/18)<1$.
Consequently $D(n)/n^{n/2}$ cannot tend to $1$ as $n$ tends to infinity
through all positive integers. This disproves the second clause, and hence
the conjunction, in the original conjecture.
\end{abstract}

\section{The original assertion and the result}

Both language versions of conjecture 00000000236 define the same maximum
$D(n)$ and assert the conjunction of the following statements:
\begin{enumerate}
\item For every positive integer $n$ divisible by $4$, there is a real
Hadamard matrix of order $n$.
\item The real sequence $D(n)/n^{n/2}$ tends to $1$ as $n\to\infty$.
\end{enumerate}
A real Hadamard matrix has entries in $\{-1,1\}$ and satisfies
$HH^{\mathsf T}=nI_n$. The divisibility condition in the first clause
does not restrict the sequence in the second clause. The exponent $n/2$
is real division, including at odd orders. The class defining $D(n)$ is
finite and nonempty, so its maximum exists and is attained.

\begin{theorem}\label{thm:gap}
For every odd integer $n\geq 3$ and every real $n\times n$ matrix $A$
with entries in $\{-1,1\}$,
\[
 \frac{(\det A)^2}{n^n}\leq e^{-1/18}.
\]
In particular,
\[
 0\leq \frac{D(n)}{n^{n/2}}\leq e^{-1/36}<1
 \qquad(n\geq3\text{ odd}).
\]
Thus the second clause is false. This argument makes no claim that the
Hadamard existence assertion in the first clause is false.
\end{theorem}

The proof below uses only parity of integer scalar products, the spectral
theorem for a real symmetric matrix, and an elementary product estimate.
Its constant is sufficient to separate the sequence from $1$ along
infinitely many orders; no optimal determinant formula is required.

\section{An elementary product estimate}

\begin{lemma}\label{lem:product}
Suppose $x_1,\ldots,x_n$ are nonnegative real numbers satisfying
\[
 \sum_{i=1}^n x_i=n,
 \qquad \sum_{i=1}^n(x_i-1)^2\geq\frac12.
\]
Then $\prod_{i=1}^n x_i\leq e^{-1/18}$.
\end{lemma}

\begin{proof}
Put $S=\sum_i(\sqrt{x_i}-1)^2$. We first show that $S\geq1/18$.
If $S<1/18$, every summand is at most $S$, and therefore
$0\leq\sqrt{x_i}<2$. It follows that
\[
 (x_i-1)^2=(\sqrt{x_i}-1)^2(\sqrt{x_i}+1)^2
 \leq9(\sqrt{x_i}-1)^2.
\]
Summation would give $1/2\leq9S<1/2$, a contradiction.

The inequality $1+u\leq e^u$, applied to $u=\sqrt{x_i}-1$, gives
$\sqrt{x_i}\leq e^{\sqrt{x_i}-1}$. Both sides are nonnegative, so
\[
 x_i\leq e^{2(\sqrt{x_i}-1)}.
\]
Moreover, using $\sum_i x_i=n$,
\[
 \sum_i2(\sqrt{x_i}-1)
 =\sum_i\bigl(x_i-1-(\sqrt{x_i}-1)^2\bigr)=-S.
\]
Multiplying the pointwise bounds yields
\[
 \prod_i x_i\leq e^{-S}\leq e^{-1/18}.
\]
This argument also covers zero values among the $x_i$.
\end{proof}

\section{Parity forces the spectral gap}

Let $n\geq3$ be odd and let $A$ be a sign matrix of order $n$.
Its Gram matrix $G=AA^{\mathsf T}$ is real symmetric and positive
semidefinite. Every diagonal entry is $n$. Every entry $G_{ij}$ is an
integer sum of $n$ signs, hence is odd and nonzero. Thus
$G_{ij}^2\geq1$ for all $i,j$, while $G_{ii}^2=n^2$.
Consequently,
\begin{align}
 \tr G&=n^2,\label{eq:trace}\\
 \tr(G^2)&=\sum_{i,j}G_{ij}^2
 \geq n^3+n(n-1).\label{eq:square}
\end{align}
The equality in \eqref{eq:square} uses the symmetry of $G$.

Write $\lambda_1,\ldots,\lambda_n\geq0$ for the real eigenvalues
of $G$, counted with multiplicity, and set $x_i=\lambda_i/n$.
The spectral theorem gives
\[
 \sum_i\lambda_i=\tr G,\qquad
 \sum_i\lambda_i^2=\tr(G^2),\qquad
 \prod_i\lambda_i=\det G=(\det A)^2.
\]
In particular $\sum_i x_i=n$ by \eqref{eq:trace}, and
\begin{align*}
 \sum_i(x_i-1)^2
 &=\frac{\tr(G^2)}{n^2}-n\\
 &\geq 1-\frac1n\geq\frac12.
\end{align*}
Lemma~\ref{lem:product} now implies
\[
 \frac{(\det A)^2}{n^n}
 =\prod_i x_i\leq e^{-1/18},
\]
proving the matrix estimate in Theorem~\ref{thm:gap}.
Choosing a matrix attaining $D(n)$ gives the same estimate for $D(n)^2$.
Since $n>0$ and $(n^{n/2})^2=n^n$, taking nonnegative square roots
gives the stated bound on $D(n)/n^{n/2}$.

\section{Failure of the original limit}

For completeness, put $c=e^{-1/36}<1$ and
$\varepsilon=(1-c)/2>0$. If $D(n)/n^{n/2}\to1$ through all positive
integers, there is a threshold $N$ beyond which
\[
 \frac{D(n)}{n^{n/2}}>1-\varepsilon
 =\frac{1+c}{2}>c.
\]
There is an odd integer $n\geq\max(N,3)$, contradicting
Theorem~\ref{thm:gap}. Therefore the asserted limit is false, and so is
the original conjunction. This is an infinite-order obstruction, with a
fixed positive gap, rather than evidence from a finite sample of matrices.

\section{Formal correspondence and verification}

The Lean project uses Lean 4.19.0 and Mathlib v4.19.0. Its namespace is
\texttt{Conjecture236}. The source-to-formalization correspondence is as
follows:
\begin{itemize}\raggedright
\item \texttt{SignMatrix n} is the finite type of Boolean matrices.
\texttt{integerMatrix} sends its entries to the integers $1$ and $-1$,
and \texttt{realMatrix} casts those entries to $\R$.
\texttt{sign\_encoding\_complete} proves that every arbitrary real sign
matrix occurs in this encoding.
\item \texttt{D} is the finite maximum of absolute determinants over the
whole encoded class. The theorem \texttt{D\_isMaximum} proves both
attainment by a real sign matrix and domination of every real sign
matrix's absolute determinant. Thus the finite representation preserves
the source's maximum.
\item \texttt{IsHadamard} uses real sign entries and the actual matrix
equation $HH^{\mathsf T}=nI_n$.
\texttt{isHadamard\_iff\_rows} identifies that equation with the usual
row scalar products. \texttt{HadamardConjecture} quantifies over all
positive natural orders divisible by four.
\item \texttt{normalizedMaximum} uses real exponentiation with exponent
$(n:\R)/2$. \texttt{normalization\_pos} and
\texttt{normalization\_square} verify the positive denominator and its
exact square for positive orders.
\item \texttt{AllOrdersLimit} is ordinary filter convergence of this
sequence to $1$. The theorem
\path{allOrdersLimit_iff_positive_epsilon} proves its equivalence
to the positive-index $\varepsilon$--$N$ formulation. A natural-number
index at zero is merely a finite extension and is never used to refute
the limit.
\end{itemize}

The spectral and parity estimates culminate in
\texttt{sign\_determinant\_gap}; the corresponding estimate for the true
maximum is \texttt{normalizedMaximum\_odd\_gap}.
\texttt{not\_AllOrdersLimit} contradicts convergence by choosing the odd
order $2N+3$ beyond an arbitrary eventual threshold $N$.
The definition \texttt{OriginalConjecture} is the conjunction of
\texttt{HadamardConjecture} and \texttt{AllOrdersLimit}. The final theorem
\texttt{conjecture236\_false} proves $\neg\,\texttt{OriginalConjecture}$.

\paragraph{Verification.}
The frozen project was rebuilt in fresh directories and its Lean sources
were replayed with warnings treated as errors, both by the coordinator
and by a separate nonauthor reviewer. The complete compiled-module audit
covers all 33 written declarations (11 definitions and 22 theorems) and
all 17 generated declarations. Their actual types, axiom dependencies,
and transitive logical dependencies were checked. Only Lean's standard
\texttt{propext}, \texttt{Classical.choice}, and \texttt{Quot.sound}
occur; there are no admitted proofs, additional axioms, or unsafe or
partial logical dependencies. The formalization uses no
\texttt{native\_decide} and assumes no external numerical computation.

The exact compiler and all nine dependency revisions were checked.
Cached library build artifacts were reused; the checks rebuild this
submission, not the entire compiler and Mathlib from source.
The source, pinned project, complete execution records, independent
review, and reproduction instructions accompany this report. These are
local verification results, distinct from acceptance by the repository's
maintainers.

\end{document}
